\documentclass[11pt,a4paper]{ctexart}
\usepackage[a4paper,margin=1.65cm]{geometry}
\usepackage{amsmath,amssymb,bm}
\usepackage{booktabs,longtable,array,multicol}
\usepackage{enumitem}
\usepackage{xcolor}
\usepackage{hyperref}
\usepackage{fancyhdr}
\usepackage{titlesec}
\hypersetup{colorlinks=true,linkcolor=blue!50!black,urlcolor=blue!50!black}
\setlist{nosep,leftmargin=1.8em}
\setlength{\parindent}{0pt}
\setlength{\parskip}{4pt}
\setlength{\headheight}{14pt}
\pagestyle{fancy}
\fancyhf{}
\lhead{工程电磁场期末复习}
\rhead{EE208}
\cfoot{\thepage}
\titleformat{\section}{\Large\bfseries\color{blue!45!black}}{\thesection}{0.6em}{}
\titleformat{\subsection}{\bfseries\color{blue!35!black}}{\thesubsection}{0.5em}{}
\newcommand{\dd}{\mathrm{d}}
\newcommand{\uv}[1]{\hat{\bm a}_{#1}}
\newcommand{\E}{\bm E}
\newcommand{\D}{\bm D}
\newcommand{\B}{\bm B}
\newcommand{\Hf}{\bm H}
\newcommand{\J}{\bm J}
\newcommand{\K}{\bm K}
\newcommand{\M}{\bm M}
\newcommand{\A}{\bm A}
\newcommand{\V}{\bm V}
\newcommand{\curl}{\nabla\times}
\newcommand{\divg}{\nabla\cdot}
\newcommand{\grad}{\nabla}

\begin{document}

\begin{center}
{\huge\bfseries 工程电磁场期末复习}\\[4pt]
{\large 基于本目录 lecture、作业与已有复习资料整理}\\[2pt]
{\small 生成日期：2026-06-03}
\end{center}

\tableofcontents
\newpage

\section{考前总览}
\textbf{主线：}源量 $\rho,\J$ 产生场 $\E,\D,\Hf,\B$，场满足积分定律和微分定律；遇到边界、对称性、能量或电路等效时，换成势函数、电容、电感、磁路或麦克斯韦方程求解。

\begin{longtable}{p{0.18\linewidth}p{0.38\linewidth}p{0.36\linewidth}}
\toprule
模块 & 必会公式 & 常见题型 \\
\midrule
矢量分析 & $\grad V,\ \divg \A,\ \curl \A$；高斯/斯托克斯定理 & 坐标系微元、通量、环量、散度/旋度判断 \\
静电场 & 库仑定律、$\D=\epsilon\E$、$\divg\D=\rho_v$、$\curl\E=0$ & 点/线/面/体电荷场，电势差，导体边界 \\
导体与介质 & $\J=\sigma\E$，连续性方程，边界条件 & 电介质交界面、电阻、弛豫时间 \\
电容与边值 & $C=Q/V$，$\nabla^2 V=-\rho_v/\epsilon$ & 平板、同轴、球形、泊松/拉普拉斯方程 \\
稳恒磁场 & 毕奥-萨伐尔、安培环路、$\B=\mu\Hf$ & 长直线、同轴线、螺线管、磁场力 \\
磁介质与磁路 & $\B=\mu_0(\Hf+\M)$，$\mathcal R=l/(\mu S)$ & 磁路、气隙、边界条件 \\
电感与时变场 & $L=\lambda/I$，$\curl\E=-\partial \B/\partial t$，$\curl\Hf=\J+\partial\D/\partial t$ & 法拉第感应、位移电流、麦克斯韦方程组 \\
\bottomrule
\end{longtable}

\textbf{做题优先级：}
\begin{enumerate}
\item 先看对称性：球对称用球面，高斯柱对称用圆柱面，平面对称用 pillbox。
\item 再选定律：闭合面积分求通量，闭合线积分求环量，势函数适合边值。
\item 最后检查单位和方向：$\E$ 为 V/m，$\D$ 为 C/m$^2$，$\Hf$ 为 A/m，$\B$ 为 T。
\end{enumerate}

\section{矢量分析速查}
\subsection{三种正交坐标微元}
\begin{longtable}{p{0.18\linewidth}p{0.25\linewidth}p{0.25\linewidth}p{0.22\linewidth}}
\toprule
坐标 & 线元 & 面元示例 & 体元 \\
\midrule
直角 $(x,y,z)$ &
$\dd \bm l=\uv{x}\dd x+\uv{y}\dd y+\uv{z}\dd z$ &
$\dd \bm S_x=\uv{x}\dd y\dd z$ &
$\dd v=\dd x\dd y\dd z$ \\
圆柱 $(\rho,\phi,z)$ &
$\dd \bm l=\uv{\rho}\dd\rho+\uv{\phi}\rho\dd\phi+\uv{z}\dd z$ &
$\dd \bm S_\rho=\uv{\rho}\rho\dd\phi\dd z$ &
$\dd v=\rho\dd\rho\dd\phi\dd z$ \\
球 $(r,\theta,\phi)$ &
$\dd \bm l=\uv{r}\dd r+\uv{\theta}r\dd\theta+\uv{\phi}r\sin\theta\dd\phi$ &
$\dd \bm S_r=\uv{r}r^2\sin\theta\dd\theta\dd\phi$ &
$\dd v=r^2\sin\theta\dd r\dd\theta\dd\phi$ \\
\bottomrule
\end{longtable}

\subsection{梯度、散度、旋度的物理意义}
\begin{itemize}
\item $\grad V$：标量场变化最快的方向；静电中 $\E=-\grad V$。
\item $\divg \A$：单位体积净流出量；$\divg\D=\rho_v$ 表示电荷是电通量源。
\item $\curl \A$：单位面积环量密度；静电场 $\curl\E=0$，稳恒磁场 $\curl\Hf=\J$。
\end{itemize}

\textbf{两个积分定理：}
\[
\oint_S \A\cdot \dd\bm S=\int_V \divg\A\,\dd v,\qquad
\oint_C \A\cdot \dd\bm l=\int_S(\curl\A)\cdot\dd\bm S .
\]

\section{静电场}
\subsection{库仑定律与电场强度}
点电荷 $Q$ 在自由空间产生
\[
\E=\frac{Q}{4\pi\epsilon_0R^2}\uv{R},\qquad
\bm F=Q_t\E .
\]
连续电荷分布：
\[
\E=\frac{1}{4\pi\epsilon}\int \frac{\rho_l\,\dd l\ \uv R}{R^2},\quad
\E=\frac{1}{4\pi\epsilon}\int \frac{\rho_s\,\dd S\ \uv R}{R^2},\quad
\E=\frac{1}{4\pi\epsilon}\int \frac{\rho_v\,\dd v\ \uv R}{R^2}.
\]

\subsection{高斯定律}
\[
\oint_S \D\cdot\dd\bm S=Q_{\rm enc},\qquad \divg\D=\rho_v,\qquad \D=\epsilon\E .
\]
高斯面选择口诀：
\begin{itemize}
\item 点电荷/球壳：球面，$\E=Q/(4\pi\epsilon r^2)\uv r$。
\item 无限长线电荷：圆柱面，$\E=\rho_l/(2\pi\epsilon\rho)\uv\rho$。
\item 无限大面电荷：pillbox，单面 $\E=\rho_s/(2\epsilon)\uv n$；导体表面外侧 $\D_n=\rho_s$。
\end{itemize}

\subsection{电势与保守场}
\[
V_{ab}=V_a-V_b=-\int_b^a \E\cdot\dd\bm l,\qquad
\E=-\grad V,\qquad \oint_C \E\cdot\dd\bm l=0 .
\]
点电荷电势：$V=Q/(4\pi\epsilon R)$，零电位通常取在无穷远。线电荷电势差常写为
\[
V_{\rho=a}-V_{\rho=b}=\frac{\rho_l}{2\pi\epsilon}\ln\frac{b}{a}.
\]

\subsection{电偶极子}
电偶极矩 $\bm p=Q\bm d$。远区近似：
\[
V\approx \frac{\bm p\cdot \uv r}{4\pi\epsilon r^2}
=\frac{p\cos\theta}{4\pi\epsilon r^2}.
\]
放在外电场中受力矩 $\bm T=\bm p\times\E$。

\section{导体、电流与电介质}
\subsection{电流密度与连续性}
\[
I=\int_S \J\cdot\dd\bm S,\qquad
\J=\rho_v\bm v,\qquad
\J=\sigma\E.
\]
电荷守恒：
\[
\divg\J=-\frac{\partial \rho_v}{\partial t}.
\]
稳恒电流时 $\divg\J=0$。良导体内部静电平衡时 $\E=0$，净电荷只分布在表面，导体表面是等势面。

\subsection{介质极化}
\[
\D=\epsilon_0\E+\bm P=\epsilon\E,\qquad
\epsilon=\epsilon_r\epsilon_0,\qquad
\bm P=\epsilon_0\chi_e\E .
\]
各向同性线性介质中只需用 $\epsilon$ 替换自由空间介电常数。

\subsection{静电边界条件}
设 $\uv n$ 从介质 1 指向介质 2：
\[
\uv n\cdot(\D_2-\D_1)=\rho_s,\qquad
\uv n\times(\E_2-\E_1)=0.
\]
若无自由面电荷，$D_{1n}=D_{2n}$；切向电场总连续，$E_{1t}=E_{2t}$。导体表面外侧满足 $E_t=0,\ D_n=\rho_s$。

\section{电容、能量与边值问题}
\subsection{电容定义和常见模型}
\[
C=\frac{Q}{V_0},\qquad W_e=\frac12CV^2=\frac12QV=\frac{Q^2}{2C},\qquad
w_e=\frac12\E\cdot\D.
\]
\begin{longtable}{p{0.25\linewidth}p{0.32\linewidth}p{0.35\linewidth}}
\toprule
结构 & 电容 & 适用条件 \\
\midrule
平行板 & $C=\epsilon S/d$ & 忽略边缘效应 \\
同轴线，单位长度 & $C'=\dfrac{2\pi\epsilon}{\ln(b/a)}$ & 内半径 $a$，外半径 $b$ \\
球形电容 & $C=4\pi\epsilon\dfrac{ab}{b-a}$ & 两同心球半径 $a,b$ \\
孤立导体球 & $C=4\pi\epsilon a$ & $b\to\infty$ \\
双线传输线，单位长度 & $C'=\dfrac{\pi\epsilon}{\operatorname{arccosh}(D/2a)}$ & 两线半径 $a$，中心距 $D$ \\
\bottomrule
\end{longtable}

\subsection{多层介质的等效}
\begin{itemize}
\item 电场垂直穿过多层介质：各层 $D_n$ 相同，等效为串联，
$V=\sum_i D d_i/\epsilon_i$。
\item 电场平行分布在不同介质区域：各区 $E_t$ 相同，等效为并联，
$Q=\sum_i \epsilon_i E S_i$。
\end{itemize}

\subsection{泊松方程与拉普拉斯方程}
由 $\divg\D=\rho_v$ 和 $\E=-\grad V$：
\[
\nabla^2V=-\frac{\rho_v}{\epsilon}.
\]
无体电荷区为拉普拉斯方程 $\nabla^2V=0$。常用一维解：
\[
\frac{\dd^2V}{\dd x^2}=0\Rightarrow V=Ax+B,\qquad
\frac{1}{\rho}\frac{\dd}{\dd\rho}\left(\rho\frac{\dd V}{\dd\rho}\right)=0\Rightarrow V=A\ln\rho+B,
\]
\[
\frac{1}{r^2}\frac{\dd}{\dd r}\left(r^2\frac{\dd V}{\dd r}\right)=0\Rightarrow V=A+\frac{B}{r}.
\]
解边值题步骤：选坐标 $\to$ 写拉普拉斯/泊松方程 $\to$ 按对称性降维 $\to$ 积分得通解 $\to$ 代边界条件 $\to$ 用 $\E=-\grad V$、$\D=\epsilon\E$、$Q=\int\D\cdot\dd\bm S$ 求目标量。

\section{稳恒磁场}
\subsection{毕奥-萨伐尔定律}
\[
\dd\Hf=\frac{I\dd\bm l\times\uv R}{4\pi R^2},\qquad
\Hf=\int \frac{I\dd\bm l\times\uv R}{4\pi R^2}.
\]
常见结果：
\[
\Hf_{\rm long\ wire}=\frac{I}{2\pi\rho}\uv\phi,\qquad
\Hf_{\rm circular\ center}=\frac{I}{2a}\uv z.
\]

\subsection{安培环路定律}
\[
\oint_C \Hf\cdot\dd\bm l=I_{\rm enc},\qquad
\curl\Hf=\J.
\]
适合高度对称的电流分布。长直导线、同轴电缆、无限大电流片、理想螺线管都优先用安培环路。

\subsection{磁通密度与磁通}
\[
\B=\mu\Hf,\qquad \Phi=\int_S\B\cdot\dd\bm S,\qquad \divg\B=0.
\]
磁场没有孤立磁荷，所以闭合曲面的净磁通为零。

\subsection{磁场力}
洛伦兹力：
\[
\bm F=Q(\E+\bm v\times\B).
\]
载流导线受力：
\[
\dd\bm F=I\dd\bm l\times\B,\qquad \bm F=\int I\dd\bm l\times\B.
\]
两无限长平行导线单位长度受力：
\[
\frac{F}{L}=\frac{\mu I_1I_2}{2\pi d}.
\]
同向相吸，反向相斥。

\section{磁介质、磁路与电感}
\subsection{磁化与边界条件}
\[
\B=\mu_0(\Hf+\M),\qquad \B=\mu\Hf\quad\text{(线性介质)}.
\]
磁边界条件：
\[
\uv n\cdot(\B_2-\B_1)=0,\qquad
\uv n\times(\Hf_2-\Hf_1)=\K.
\]
若无表面电流，$B_n$ 连续、$H_t$ 连续。

\subsection{磁路}
\[
\mathcal F=NI,\qquad \mathcal R=\frac{l}{\mu S},\qquad
\Phi=\frac{\mathcal F}{\mathcal R}.
\]
气隙磁阻常远大于铁芯磁阻，磁路题先算各段 $\mathcal R$ 串并联，再由 $\Phi$ 得 $B=\Phi/S$，由 $\Hf=B/\mu$ 得各段磁场。

\subsection{电感}
\[
L=\frac{\lambda}{I}=\frac{N\Phi}{I},\qquad
W_m=\frac12LI^2,\qquad
w_m=\frac12\B\cdot\Hf.
\]
同轴线单位长度电感：
\[
L'=\frac{\mu}{2\pi}\ln\frac{b}{a}.
\]
长螺线管近似：
\[
L=\frac{\mu N^2S}{l}.
\]

\section{时变场与麦克斯韦方程}
\subsection{法拉第电磁感应}
\[
\oint_C\E\cdot\dd\bm l=-\frac{\dd}{\dd t}\int_S\B\cdot\dd\bm S,\qquad
\curl\E=-\frac{\partial\B}{\partial t}.
\]
若回路运动，感应电动势可写为
\[
\mathcal E=\oint_C(\E+\bm v\times\B)\cdot\dd\bm l.
\]
方向用楞次定律：感应电流产生的磁场总是反抗磁通变化。

\subsection{位移电流}
安培定律在时变场中修正为
\[
\oint_C\Hf\cdot\dd\bm l=\int_S\left(\J+\frac{\partial\D}{\partial t}\right)\cdot\dd\bm S,
\qquad
\curl\Hf=\J+\frac{\partial\D}{\partial t}.
\]
电容器极板间没有传导电流，但有位移电流密度 $\partial\D/\partial t$，保证电流连续。

\subsection{麦克斯韦方程组}
\begin{longtable}{p{0.46\linewidth}p{0.46\linewidth}}
\toprule
积分形式 & 微分形式 \\
\midrule
$\displaystyle \oint_S\D\cdot\dd\bm S=Q_{\rm enc}$ &
$\displaystyle \divg\D=\rho_v$ \\
$\displaystyle \oint_S\B\cdot\dd\bm S=0$ &
$\displaystyle \divg\B=0$ \\
$\displaystyle \oint_C\E\cdot\dd\bm l=-\dfrac{\dd}{\dd t}\int_S\B\cdot\dd\bm S$ &
$\displaystyle \curl\E=-\dfrac{\partial\B}{\partial t}$ \\
$\displaystyle \oint_C\Hf\cdot\dd\bm l=\int_S(\J+\partial\D/\partial t)\cdot\dd\bm S$ &
$\displaystyle \curl\Hf=\J+\dfrac{\partial\D}{\partial t}$ \\
\bottomrule
\end{longtable}

\section{典型题型模板}
\subsection{由电荷分布求电场}
\begin{enumerate}
\item 判断是否能用高斯定律。若场在高斯面上大小恒定且方向与面元平行/垂直，优先高斯。
\item 否则写库仑积分，确定源点、场点、$\bm R$ 和 $\uv R$。
\item 利用对称性消去分量，积分剩余分量。
\item 检查远场极限，例如有限带电体远处应近似点电荷场。
\end{enumerate}

\subsection{电容题}
\begin{enumerate}
\item 假设导体带 $+Q,-Q$ 或单位长度电荷 $\rho_l$。
\item 用高斯定律求 $\D$ 和 $\E$。
\item 积分求电势差 $V=\int \E\cdot\dd\bm l$。
\item 代入 $C=Q/V$ 或 $C'=\rho_l/V$。
\end{enumerate}

\subsection{边界条件题}
\begin{enumerate}
\item 把场分解为法向和切向：$\E=\E_n+\E_t$，$\D=\D_n+\D_t$。
\item 静电边界：$D_{2n}-D_{1n}=\rho_s$，$E_{2t}=E_{1t}$。
\item 磁场边界：$B_{2n}=B_{1n}$，$\Hf_{2t}-\Hf_{1t}=\K\times\uv n$ 的等价方向式。
\item 注意 $\D=\epsilon\E$、$\B=\mu\Hf$ 会让折射角改变。
\end{enumerate}

\subsection{磁场题}
\begin{enumerate}
\item 长直、同轴、螺线管：选安培环路，$\oint\Hf\cdot\dd\bm l=I_{\rm enc}$。
\item 任意线电流或圆环：用毕奥-萨伐尔。
\item 求力时先求外磁场，再用 $\dd\bm F=I\dd\bm l\times\B$。
\item 求电感时先求 $\Phi$ 或磁能，再用 $L=N\Phi/I$ 或 $W_m=\frac12LI^2$。
\end{enumerate}

\subsection{法拉第感应题}
\begin{enumerate}
\item 明确积分面方向和回路正方向，二者满足右手定则。
\item 计算磁通 $\Phi=\int\B\cdot\dd\bm S$。
\item 用 $\mathcal E=-\dd\Phi/\dd t$，运动导体可用 $\oint(\bm v\times\B)\cdot\dd\bm l$。
\item 用楞次定律核对电流方向。
\end{enumerate}

\section{最后一页：高频公式清单}
\begin{multicols}{2}
\begin{align*}
&\E=\frac{Q}{4\pi\epsilon R^2}\uv R\\
&\D=\epsilon\E,\quad \B=\mu\Hf\\
&\oint_S\D\cdot\dd\bm S=Q\\
&\oint_C\E\cdot\dd\bm l=0\quad(\text{静电})\\
&V_{ab}=-\int_b^a\E\cdot\dd\bm l\\
&\E=-\grad V\\
&\J=\sigma\E,\quad I=\int_S\J\cdot\dd\bm S\\
&\divg\J=-\partial\rho_v/\partial t\\
&C=Q/V,\quad W_e=\tfrac12CV^2\\
&\nabla^2V=-\rho_v/\epsilon\\
&\dd\Hf=\frac{I\dd\bm l\times\uv R}{4\pi R^2}\\
&\oint_C\Hf\cdot\dd\bm l=I\\
&\curl\Hf=\J,\quad \divg\B=0\\
&\bm F=Q(\E+\bm v\times\B)\\
&\dd\bm F=I\dd\bm l\times\B\\
&\mathcal R=l/(\mu S),\quad \Phi=NI/\mathcal R\\
&L=N\Phi/I,\quad W_m=\tfrac12LI^2\\
&\curl\E=-\partial\B/\partial t\\
&\curl\Hf=\J+\partial\D/\partial t
\end{align*}
\end{multicols}

\textbf{考场提醒：}不要急着代公式。先画图、标方向、写单位，再决定积分路径或积分面。多数题的关键不是积分难，而是高斯面、安培环路、边界方向和参考电位选得是否干净。

\end{document}
